% 完整翻译：原PDF物理页63--82，原书页56--75。
\SourcePage{63}
\section*{2.1. 分层空间}
\addcontentsline{toc}{section}{2.1. 分层空间}

\par\medskip\noindent\textbf{2.1.1.} 设 $X$ 是带有环\term{sheaf}{层} $\mathcal O$ 的拓扑空间，$\mathcal S$ 是将 $X$ 分成有限个局部闭子集的一个划分（这些子集称为\term{stratum}{层片}），$p\colon\mathcal S\to\Z$ 是一个函数（称为\term{perversity}{反常度}）。按“划分”的定义，各\term{stratum}{层片}均非空。假定每个\term{stratum}{层片}的闭包都是若干\term{stratum}{层片}的并。这些条件可以放宽，见3.2.19。

回顾0.0：对于 $f\colon X\to Y$，我们将通常记为 $Rf_!,Rf_*,Rf^!,Rf^*$（或 $Lf^*$）的\term{derived-category}{导出范畴}之间的函子，简记为 $f_!,f_*,f^!,f^*$。由于2.1.7中说明的原因，\term{sheaf}{层}范畴之间的同名函子将在左上角加 $0$。

\par\medskip\noindent\textbf{定义2.1.2.} $\D(X,\mathcal O)$ 的子范畴 ${}^p\D^{\leq0}(X,\mathcal O)$（分别地，${}^p\D^{\geq0}(X,\mathcal O)$）由满足下述条件的\term{complex}{复形} $K$（分别地，$K\in\D^+(X,\mathcal O)$）组成：对每个\term{stratum}{层片} $S$，记其到 $X$ 的嵌入为 $i_S$，则
\[
 H^ni_S^*K=0\quad(n>p(S)),\qquad
 \text{分别地，}\quad H^ni_S^!K=0\quad(n<p(S)).
\]

函子 ${}^0i_S^*$ 的正合性使我们可以把 ${}^p\D^{\leq0}(X,\mathcal O)$ 的定义改述为：$K\in{}^p\D^{\leq0}(X,\mathcal O)$ 当且仅当 $H^iK$ 在 $S$ 上的限制对 $i>p(S)$ 为零。因此，自然\term{t-structure}{$t$-结构}的函子 $\tau_{\leq a}$ 与 $\tau_{\geq a}$ 都把 ${}^p\D^{\leq0}(X,\mathcal O)$ 映入自身。

若各函子 ${}^0i_S^!$ 的\term{cohomology}{上同调}维数有限，则函子 $i_S^!\colon\D^+(X,\mathcal O)\to\D^+(S,\mathcal O)$ 自然延伸为 $\D(X,\mathcal O)\to\D(S,\mathcal O)$。于是，即使 $K$ 不一定属于 $\D^+(X,\mathcal O)$，条件“$H^ni_S^!K=0$，$n<p(S)$”仍有意义。不过，这个条件蕴含 $K\in\D^+(X,\mathcal O)$。更精确地说，若 $p\geq a$，它就蕴含 $K\in\D^{\geq a}(X,\mathcal O)$。我们按照 $S_1\subset\overline{S_2}$ 所给出的\term{stratum}{层片}次序作下降归纳，证明 $H^iK$ 在 $S$ 上对 $i<a$ 为零。由三角 $(\tau_{<a}K,K,\tau_{\geq a}K)$ 以及 ${}^0i_S^!$ 的左正合性，对 $i<a$ 有
\[
 H^ii_S^!\tau_{<a}K\xrightarrow{\sim}H^ii_S^!K.
\]
归纳假设保证，$H^j\tau_{<a}K$ 在满足 $T\ne S$ 且 $S\subset\overline T$ 的\term{stratum}{层片} $T$ 上为零。因此，$S$ 有一个邻域，在这个邻域内 $H^j\tau_{<a}K$ 的\term{support}{支撑}包含于 $S$。故对 $i<a$，
\[
 H^iK|_S=H^ii_S^*\tau_{<a}K
 =H^ii_S^!\tau_{<a}K=H^ii_S^!K,
\]
最后一项按假设为零。

\SourcePage{64}
不作\term{cohomology}{上同调}维数的假设时，同一论证仍适用于 $K\in\D^+(X,\mathcal O)$。若整数 $a,b$ 满足 $a\leq p\leq b$，则
\begin{equation*}
 \D^{\leq a}(X,\mathcal O)\subset{}^p\D^{\leq0}(X,\mathcal O)
 \subset\D^{\leq b}(X,\mathcal O),\tag{2.1.2.1}
\end{equation*}
以及
\[
 \D^{\geq a}(X,\mathcal O)\supset{}^p\D^{\geq0}(X,\mathcal O)
 \supset\D^{\geq b}(X,\mathcal O).
\]
用 ${}^p\D^{+,\leq0}(X,\mathcal O)$ 表示 $\D^+(X,\mathcal O)$ 与 ${}^p\D^{\leq0}(X,\mathcal O)$ 的交。把 $+$ 换成 $-$ 或 $b$、把 $0$ 换成 $n$ 时，也使用同样的记法。

\par\medskip\noindent\textbf{命题2.1.3.} 对每个\term{perversity}{反常度} $p$，$({}^p\D^{+,\leq0}(X,\mathcal O),{}^p\D^{\geq0}(X,\mathcal O))$ 是 $\D^+(X,\mathcal O)$ 上的一个 \term{t-structure}{$t$-结构}。

对\term{stratum}{层片}数 $N$ 作归纳。$N=0$ 时 $X=\varnothing$，断言显然成立。$N=1$ 时，得到的就是自然 \term{t-structure}{$t$-结构}平移 $p(X)$ 所得的结构。设 $N\geq2$，取 $X$ 的一个真闭子集 $F$，使它是\term{stratum}{层片}的并，并令 $U$ 为其补开集。例如，可以取 $F$ 为一个闭\term{stratum}{层片}。归纳假设适用于带有诱导\term{stratification}{分层}的 $F$ 与 $U$，给出 $\D^+(U,\mathcal O)$ 和 $\D^+(F,\mathcal O)$ 上的 \term{t-structure}{$t$-结构}。所求的 $\D^+(X,\mathcal O)$ 上的 \term{t-structure}{$t$-结构}由1.4.10得到。

\par\medskip\noindent\textbf{推论2.1.4.} $({}^p\D^{\leq0}(X,\mathcal O),{}^p\D^{\geq0}(X,\mathcal O))$ 是 $\D(X,\mathcal O)$ 上的一个 \term{t-structure}{$t$-结构}。它在 $\D^*(X,\mathcal O)$ 上诱导 \term{t-structure}{$t$-结构}，其中 $*=+,-,b$。

取整数 $a,b$ 使 $a\leq p\leq b$。证明中固定 $(X,\mathcal O)$，将它从记号中略去。设 $K\in{}^p\D^{\leq0}$、$L\in{}^p\D^{>0}$，来验证 $\Hom(K,L)=0$。由于 $L\in\D^{>a}$，有 $\Hom(\tau_{\leq a}K,L)=0$；又因 $\tau_{>a}K\in{}^p\D^{+,\leq0}$，可应用2.1.3，得到 $\Hom(\tau_{>a}K,L)=0$。由\term{distinguished-triangle}{区别三角} $(\tau_{\leq a}K,K,\tau_{>a}K)$ 的 $\Hom$ 长\term{exact-sequence}{正合列}即得结论。最后，设 $K\in\D$。依2.1.3，有\term{distinguished-triangle}{区别三角} $(U,\tau_{>a}K,V)$，其中 $U\in{}^p\D^{+,\leq0}$、$V\in{}^p\D^{>0}$。应用TR4，得到两个\term{distinguished-triangle}{区别三角} $(\tau_{\leq a}K,Y,U)$ 与 $(Y,K,V)$。前者表明 $Y\in{}^p\D^{\leq0}$，后者验证了1.3.1(iii)。公理1.3.1(ii)显然成立，故得到了 $\D$ 上的一个 \term{t-structure}{$t$-结构}，称为\emph{\term{perversity}{反常度}为 $p$ 的 \term{t-structure}{$t$-结构}}。

将相应的\term{truncation}{截断}函子记为 ${}^p\tau$。由于 ${}^p\D^{\leq0}\subset\D^{\leq b}$、${}^p\D^{\geq0}\subset\D^{\geq a}$，\term{distinguished-triangle}{区别三角} $({}^p\tau_{\leq0}K,K,{}^p\tau_{\geq1}K)$ 的通常\term{cohomology}{上同调}长\term{exact-sequence}{正合列}给出：对 $i<a$，$H^i{}^p\tau_{\leq0}K=H^iK$。同样，对 $i>b$，$H^i{}^p\tau_{\geq0}K=H^iK$。由此可知，${}^p\tau_{\leq0}$ 与 ${}^p\tau_{\geq0}$
\SourcePage{65}
保持 $\D^*$（$*=+,-,b$）。

\par\medskip\noindent\textbf{定义2.1.5.} $X$ 上 $p$-\term{perverse-sheaf}{反常层}（取值于 $\mathcal O$-模）的范畴 $\mathcal M(p,X,\mathcal O)$ 定义为
\[
 {}^p\D^{\leq0}(X,\mathcal O)\cap{}^p\D^{\geq0}(X,\mathcal O).
\]
这是 $\D^b(X,\mathcal O)$ 的一个\term{admissible}{容许}\term{abelian-subcategory}{阿贝尔子范畴}。

\par\medskip\noindent\textbf{命题2.1.6.} 设 $U$ 是 $X$ 的一个局部闭子集，且是\term{stratum}{层片}的并，$j$ 是 $U$ 到 $X$ 的嵌入。对任意\term{perversity}{反常度} $p$，函子 $j_!\colon\D(U,\mathcal O)\to\D(X,\mathcal O)$ 与 $j^*\colon\D(X,\mathcal O)\to\D(U,\mathcal O)$ 均为右 \term{t-exact}{$t$-正合}。

验证直接，留给读者。由1.3.17(iii)及其补充1.3.18(i)，$j_!$ 与 $j^*$ 的右\term{adjunction}{伴随} $j^!$ 与 $j_*$ 因而为左 \term{t-exact}{$t$-正合}。

\par\medskip\noindent\textbf{2.1.7. 记号。} 从记号中略去 $(X,\mathcal O)$。有时将 ${}^p\D^{\leq0}$（分别地，${}^p\D^{\geq0}$）记作 $\D^{\leq p}$（分别地，$\D^{\geq p}$）。当 $p$ 恒等于 $a$ 时，$\D^{\leq p}=\D^{\leq a}$（右端取自然 \term{t-structure}{$t$-结构}），且 $\D^{\geq p}=\D^{\geq a}$。对任意整数 $n$，有 $\D^{\leq p+n}={}^p\D^{\leq n}$、$\D^{\geq p+n}={}^p\D^{\geq n}$。最后，若 $p\leq q$，则 $\D^{\leq p}\subset\D^{\leq q}$、$\D^{\geq p}\supset\D^{\geq q}$。这是(2.1.2.1)的推广。同样，用 $\tau_{\leq p}$ 和 $\tau_{\geq p}$ 表示 ${}^p\tau_{\leq0}$ 和 ${}^p\tau_{\geq0}$，用 $H^p$ 表示\term{perversity}{反常度}为 $p$ 的 \term{t-structure}{$t$-结构}意义下的 $H^0$。

在2.1.6的情形中，将\term{sheaf}{层}范畴之间同名函子的导出函子简记为 $j_!,j^!,j_*,j^*$；由它们诱导的 $p$-\term{perverse-sheaf}{反常层}之间的函子在左上角加 $p$。例如，对于 $A\in\mathcal M(p,U,\mathcal O)$，置
\[
 {}^pj_!A=\tau_{\geq p}j_!A=H^p(j_!A).
\]
依1.3.18(i)，$({}^pj_!,{}^pj^!)$ 与 $({}^pj^*,{}^pj_*)$ 是两对\term{adjunction}{伴随}函子。

通常\term{sheaf}{层}的函子 $j_!,j^!,j_*,j^*$ 在左上角加 $0$：它们对应于\term{perversity}{反常度} $0$。

对 $U$ 上的一个 $p$-\term{perverse-sheaf}{反常层} $A$，$j_!A\in\D^{\leq p}(X,\mathcal O)$，而 $j_*A\in\D^{\geq p}(X,\mathcal O)$。因此，自然态射 $\alpha\colon j_!A\to j_*A$ 有分解
\[
 j_!A\longrightarrow{}^pj_!A\xrightarrow{\beta}{}^pj_*A
 \longrightarrow j_*A,\qquad\beta={}^pH^0(\alpha).
\]
定义函子 ${}^pj_{!*}$（也简记为 $j_{!*}$）如下：
\[
 j_{!*}A=\im({}^pj_!A\longrightarrow{}^pj_*A).
\]

\SourcePage{66}
对 $X$ 上的 $p$-\term{perverse-sheaf}{反常层} $A$，定义典范态射 ${}^pj^!A\to{}^pj^*A$ 为复合
\[
 {}^pj^!A\longrightarrow j^!A\longrightarrow j^*A\longrightarrow{}^pj^*A.
\]
设 $U\xrightarrow{k}V\xrightarrow{j}X$ 是由\term{stratum}{层片}组成的局部闭子集。传递公式
\[
 (jk)_!=j_!k_!,\quad(jk)_*=j_*k_*,\quad
 (jk)^!=k^!j^!,\quad(jk)^*=k^*j^*
\]
由1.3.17(iv)给出 $p$-\term{perverse-sheaf}{反常层}上的相应公式。将 ${}^pj_!$、${}^pj_{!*}$ 与 ${}^pj_*$ 作用于态射 ${}^pk_!\twoheadrightarrow{}^pk_{!*}\hookrightarrow{}^pk_*$。所得链
\[
\begin{aligned}
{}^p(jk)_!={}&{}^pj_!{}^pk_!
 \twoheadrightarrow{}^pj_!{}^pk_{!*}
 \twoheadrightarrow{}^pj_{!*}{}^pk_{!*}\\
 &\hookrightarrow{}^pj_*{}^pk_{!*}
 \hookrightarrow{}^pj_*{}^pk_*={}^p(jk)_*
\end{aligned}
\]
给出函子同构
\begin{equation*}
 {}^p(jk)_{!*}={}^pj_{!*}{}^pk_{!*}.\tag{2.1.7.1}
\end{equation*}

\par\medskip\noindent\textbf{2.1.8. 补充。} 设 $U$ 是 $X$ 中由\term{stratum}{层片}组成的开集，$F$ 为其补闭集，均赋予诱导的\term{stratification}{分层}和\term{perversity}{反常度}。于是，2.1.3给出的 $\D^+(X,\mathcal O)$ 上的 \term{t-structure}{$t$-结构}，由 $\D^+(U,\mathcal O)$ 与 $\D^+(F,\mathcal O)$ 上的 \term{t-structure}{$t$-结构}经1.4.10的方法得到。因此，1.4的形式体系可以应用。记 $F$ 到 $X$ 的嵌入为 $i$，$U$ 到 $X$ 的嵌入为 $j$。按2.1.7用 $i$ 或 $j$ 记、在上下标处加 $!$ 或 $*$、在左上角加或不加 $p$ 的函子，以及 $j_{!*}$，均与1.4中的同名函子一致。

\par\medskip\noindent\textbf{命题2.1.9.} 对 $B\in\mathcal M(p,U,\mathcal O)$，$j_{!*}B$ 是 $B$ 在 $\D(X,\mathcal O)$ 中唯一满足下列条件的延拓 $P$：对每个\term{stratum}{层片} $S\subset F$，记其到 $X$ 的嵌入为 $s$，则 $H^is^*P=0$（$i\geq p(S)$），且 $H^is^!P=0$（$i\leq p(S)$）。

更精确地说，由(1.4.14.1)，若 $\D'$ 是 $\D(X,\mathcal O)$ 中由满足上述条件的 $K$ 组成的子范畴（对每个\term{stratum}{层片} $s\colon S\to F$ 都要求这些条件成立），则 $j^*$ 诱导一个\term{equivalence}{等价} $\D'\to\D(U,\mathcal O)$。它在 $\D'\cap\mathcal M(p,X,\mathcal O)$ 上的限制给出\term{equivalence}{等价}
\[
 \D'\cap\mathcal M(p,X,\mathcal O)\longrightarrow\mathcal M(p,U,\mathcal O),
\]
其逆为\term{intermediate-extension}{中间延拓} $j_{!*}$。

${}^pj_!B$ 和 ${}^pj_*B$ 也有类似的刻画（参见1.4.23、1.3.14）。

\par\medskip\noindent\textbf{2.1.10.} 假定\term{perversity}{反常度} $p$ 满足
\begin{equation*}
 S\subset\overline T\quad\Longrightarrow\quad p(S)\geq p(T).\tag{2.1.10.1}
\end{equation*}
于是，对每个 $n$，满足 $p(S)\geq n$（分别地，$p(S)\leq n$）的\term{stratum}{层片} $S$ 之并 $F_n$（分别地，$U_n$）是闭集（分别地，开集）。记 $U_{n-1}$ 到 $U_n$ 的嵌入为 $j_n$。

\SourcePage{67}
\par\medskip\noindent\textbf{命题2.1.11.} 在2.1.10的假设和记号下，设 $A$ 是 $U_k$ 上的 $p$-\term{perverse-sheaf}{反常层}，$a$ 是满足 $a\geq k$ 且 $p\leq a$ 的整数，$j$ 为 $U_k$ 到 $X=U_a$ 的嵌入。则
\[
 j_{!*}A=\tau_{\leq a-1}j_{a*}\cdots\tau_{\leq k}j_{k+1*}A.
\]
这里的 $\tau_{\leq i}$ 均取自然 \term{t-structure}{$t$-结构}。由(2.1.7.1)，只需验证 $(j_{k+1})_{!*}A=\tau_{\leq k}j_{k+1*}A$。令 $F=U_{k+1}-U_k$。根据1.4.23，有
\[
 (j_{k+1})_{!*}A={}^p\tau^F_{\leq-1}(j_{k+1})_*A.
\]
在 $F$ 上，函数 $p$ 恒等于 $k+1$，故 ${}^p\tau^F_{\leq-1}$ 就是 $\tau^F_{\leq k}$（这里取 $F$ 上的自然 \term{t-structure}{$t$-结构}）。由于 $U_k$ 上有 $p\leq k$，由(2.1.2.1)，$A\in\D^{\leq k}(U_k,\mathcal O)$，而且
\[
 \tau_{\leq k}(j_{k+1})_*A\xrightarrow{\sim}
 \tau^F_{\leq k}(j_{k+1})_*A.
\]

\par\medskip\noindent\textbf{2.1.12.} 说明我们的记法与 M. Goresky、R. MacPherson [4]、[5] 的记法之间的关系。

(a) 他们研究以域 $R$（主要是 $R=\Q$ 或 $\C$）为系数的\term{cohomology}{上同调}，即取 $\mathcal O$ 为取值于 $R$ 的\term{constant-sheaf}{常值层}。

(b) 他们的\term{stratum}{层片}都是各处维数相同的拓扑流形，\term{stratification}{分层}满足沿各\term{stratum}{层片}的局部平凡性条件；这些条件保证：对一个\term{stratum}{层片} $S\xrightarrow{j}X$ 以及 $S$ 上的有限维 $R$-向量空间\term{local-system}{局部常值层} $F$，各 $R^ij_*F$ 在每个\term{stratum}{层片}上的限制仍局部常值且有限维。若两个\term{stratum}{层片} $S,T$ 满足 $S\subset\overline T-T$，则 $\dim S<\dim T$。

(c) $X$ 有一个开稠密\term{stratum}{层片} $U_0$，其补集的余维数满足 $\dim U_0-\dim(X-U_0)\geq2$。

(d) \term{perversity}{反常度}函数只依赖于维数：存在 $\bar p\colon\mathbb N\to\Z$，使 $p(S)=\bar p(\dim S)$。假定 $\bar p$ 递减（从而满足(2.1.10.1)）、非负，且 $p(U_0)=0$。

设 $j$ 是 $U_0$ 到 $X$ 的嵌入。核心对象是 $\D^b(X,R)$ 中的 $j_{!*}(R)$。若按2.1.11描述它，便会用到函数 $p'$：仍有 $p'(U_0)=0$，而对每个 $S\ne U_0$，令 $p'(S)=p(S)-1$。操作 $j_{!*}$ 是从 $U=U_0$ 上的\term{constant-sheaf}{常值层} $R$ 出发，依次向 $U$ 添入维数为 $\dim U_0-2,\dim U_0-3,\ldots$ 的\term{stratum}{层片} $S$；每次取\term{direct-image}{正像}，再以 $\tau_{\leq p'(S)}$ \term{truncation}{截断}。
\SourcePage{68}
他们称为“\term{perversity}{反常度}”的正是这个函数 $p'$。从本文的观点看，这相当于把 ${}^pj_{!*}$ 描述为 ${}^{p'}j_*$，因而在描述\term{duality}{对偶}现象时会出现位移。

Goresky 与 MacPherson 还假定 $\bar p$ 不会下降过快：$\bar p(n)-\bar p(n+1)=0$ 或 $1$。下面的2.1.14表明，由此可以得到对\term{stratification}{分层}选择的独立性。

\par\medskip\noindent\textbf{2.1.13.} 保留2.1.12中的下列假设：

(a) $\mathcal O$ 是取值于 $R$ 的\term{constant-sheaf}{常值层}，其中 $R$ 为左\term{noetherian}{诺特}环。

(b) 各\term{stratum}{层片}都是各处维数相同的拓扑流形。若\term{stratum}{层片} $S$ 包含于另一\term{stratum}{层片} $T$ 的闭包中，则 $\dim S<\dim T$。

(c) 对每个\term{stratum}{层片} $j\colon S\to X$，函子 ${}^0j_*$ 在 $R$-模\term{sheaf}{层}范畴上具有有限\term{cohomology}{上同调}维数。若 $F$ 是 $S$ 上有限生成 $R$-模的\term{local-system}{局部常值层}，则 $R^ij_*F$ 在每个\term{stratum}{层片}上仍局部常值且有限生成。

对于由\term{stratum}{层片}组成的局部闭子集 $U$，用 $\D_c(U,R)$（需要指明 $\mathcal S$ 时写作 $\D_{\mathcal S}(U,R)$）表示 $\D(U,R)$ 的如下全三角子范畴：其对象为\term{constructible}{可构造}的 $K$，即每个 $H^iK$ 在每个\term{stratum}{层片}上都局部常值且有限生成。同样定义 $\D_c^*(U,R)$（$*=+,-,b$）。条件(c)保证，当 $j\colon U\hookrightarrow V$ 且 $U,V$ 均为\term{stratum}{层片}组成的局部闭子集时，函子 $j_!,j^!,j_*,j^*$ 保持这些子范畴。事实上，$j_!$ 和 $j^*$ 的情形显然。$U$ 仅有一个\term{stratum}{层片}时，$j_*$ 的情形由\term{spectral-sequence}{谱序列}
\[
 R^pj_*H^qK\Longrightarrow H^{p+q}Rj_*K
\]
得出。一般情形对 $U$ 所含的\term{stratum}{层片}数归纳。若 $k\colon S\hookrightarrow U$ 是一个开\term{stratum}{层片}的嵌入，并以\term{distinguished-triangle}{区别三角} $(K,k_*k^*K,L)$ 定义 $L$，则 $L$ 是\term{constructible}{可构造}的，因为 $k_*k^*K$ \term{constructible}{可构造}；且 $L$ 的\term{support}{支撑}在 $U'=U-S$ 中。归纳假设适用于 $U'$ 到 $X$ 的嵌入，因而 $j_*L$ \term{constructible}{可构造}。再由三角 $(j_*K,(jk)_*k^*K,j_*L)$，得 $j_*K$ 也\term{constructible}{可构造}。最后，$j^!$ 的情形可归结到 $j$ 为闭嵌入的情形，再由补开集的嵌入 $k$ 对应的 $k_*$ 的情形推出：使用\term{distinguished-triangle}{区别三角} $(j_!j^!K,K,k_*k^*K)$。

若闭集 $F\subset U$ 是\term{stratum}{层片}的并，则 $\tau^F_{\leq a}$ 显然保持 $\D_c(U,R)$。由1.4.10的证明可知，对每个\term{perversity}{反常度} $p$，$\tau_{\leq p}$ 和 $\tau_{\geq p}$ 均保持 $\D_c(X,R)$。所谓 $p$-$\mathcal S$-\term{perverse-sheaf}{反常层}，就是
\SourcePage{69}
$\D_c(X,R)$ 中的 $p$-\term{perverse-sheaf}{反常层}。

设 $\mathcal T$ 是 $X$ 的一个细化 $\mathcal S$ 的\term{stratification}{分层}，并仍满足(b)、(c)。

\par\medskip\noindent\textbf{命题2.1.14.} 设 $p$ 为 $\mathcal S$ 上的\term{perversity}{反常度}，$q$ 为 $\mathcal T$ 上的\term{perversity}{反常度}。假定每当 $\mathcal S$ 的\term{stratum}{层片} $S$ 包含 $\mathcal T$ 的\term{stratum}{层片} $T$ 时，均有
\[
 p(S)\leq q(T)\leq p(S)+\dim S-\dim T.
\]
则 $\D_{\mathcal T}(X,R)$ 上\term{perversity}{反常度}为 $q$ 的 \term{t-structure}{$t$-结构}，在 $\D_{\mathcal S}(X,R)$ 上诱导\term{perversity}{反常度}为 $p$ 的 \term{t-structure}{$t$-结构}。

只需验证 $\D_{\mathcal S}^{\leq p}(X,R)\subset\D_{\mathcal T}^{\leq q}(X,R)$ 和 $\D_{\mathcal S}^{\geq p}(X,R)\subset\D_{\mathcal T}^{\geq q}(X,R)$。第一个包含关系立即由 $S\supset T$ 时 $p(S)\leq q(T)$ 得到。来验证第二个。设 $K\in\D_{\mathcal S}^{\geq p}(X,R)$，$T$ 是 $\mathcal T$ 的\term{stratum}{层片}，$S$ 是包含它的 $\mathcal S$ 的\term{stratum}{层片}，记 $T\xrightarrow{i}S\xrightarrow{j}X$。需证明 $H^n(ji)^!K=0$（$n<q(T)$）。有 $(ji)^!K=i^!j^!K$；而已知 $H^mj^!K$ 是 $S$ 上的\term{local-system}{局部常值层}，对 $m<p(S)$ 为零。记 $T$、$S$ 的定向\term{sheaf}{层}分别为 $\operatorname{or}_T$、$\operatorname{or}_S$；它们局部同构于\term{constant-sheaf}{常值层} $\Z$。设
\[
 \operatorname{or}=\mathcal Hom(\operatorname{or}_T,i^*\operatorname{or}_S)
\]
为 $T$ 在 $S$ 中的法向定向\term{sheaf}{层}，置 $d=\dim S-\dim T$。若 $L$ 是 $S$ 上\term{cohomology}{上同调}\term{sheaf}{层}局部常值的\term{sheaf}{层}\term{complex}{复形}，则已知
\[
 i^!L\simeq i^*L\otimes_{\Z}\operatorname{or}[-d].
\]
取 $L=j^!K$，得 $H^ni^!j^!K=0$（$n<p(S)+d$），而假设正是 $q(T)\leq p(S)+d$。

\par\medskip\noindent\textbf{推论2.1.15.} 每个 $p$-$\mathcal S$-\term{perverse-sheaf}{反常层}都是 $q$-$\mathcal T$-\term{perverse-sheaf}{反常层}，且包含函子正合。对由 $\mathcal S$ 的\term{stratum}{层片}组成的局部闭子集 $U$ 及其包含态射 $j\colon U\to X$，函子
\[
 {}^qj_!,{}^qj^!,{}^qj_*,{}^qj^*,{}^qj_{!*}
\]
限制到 $p$-$\mathcal S$-\term{perverse-sheaf}{反常层}上，分别就是 ${}^pj_!,{}^pj^!,{}^pj_*,{}^pj^*,{}^pj_{!*}$。

\par\medskip\noindent\textbf{2.1.16.} 除2.1.13的条件外，假定 $R$ 为交换域，$X$ 有一个局部有限的三角剖分，使 $\mathcal S$ 的每个\term{stratum}{层片} $S$ 都是开单形的并。例如，$X$ 为实代数空间，带有一个Whitney\term{stratification}{分层}。此时可应用Verdier\term{duality}{对偶}理论。Verdier\term{duality}{对偶}是 $\D_c(X,R)$ 的一个对合自同构；对于由\term{stratum}{层片}组成的局部闭子集 $j\colon U\to X$，它交换函子 $j_!$ 与 $j_*$，以及 $j^!$ 与 $j^*$。若不假定 $R$ 交换，它则把 $\D_c(X,R)$ 映到 $\D_c(X,R^{\mathrm{opp}})$。

\SourcePage{70}
对每个维数为 $d$、定向\term{sheaf}{层}为 $\operatorname{or}$ 的\term{stratum}{层片} $S$，$S$ 上的Verdier\term{duality}{对偶} $\mathbb D$（即 $K\mapsto\RHom(K,R\otimes\operatorname{or}[d])$）对 $K\in\D_S(S,R)$ 满足
\[
 H^i\mathbb DK=(H^{-d-i}K)^\vee\otimes\operatorname{or}.
\]
这里，$K$ 的\term{cohomology}{上同调}\term{sheaf}{层}局部常值且有限秩，以及 $R$ 是域，都是关键条件。（交换的局部\term{artinian}{阿廷}环也可以，例如 $\Z/\ell^n\Z$，因为它有\term{injective}{内射}的\term{duality}{对偶}化模 $I$；对 $\Z/\ell^n\Z$，可取 $I=\Z/\ell^n\Z$。此时应取的Verdier\term{duality}{对偶}为 $K\mapsto\RHom(K,Ra^!I)$，其中 $a$ 是到一点的投影。）

称下式给出的\term{perversity}{反常度} $p^*$ 为 $p$ 的\emph{\term{duality}{对偶}\term{perversity}{反常度}}：
\[
 p^*(S)=-p(S)-\dim S.
\]
上述结果表明，$\mathbb D$ 交换 $\D^{\geq p}$ 与 $\D^{\leq p^*}$，也交换 $\D^{\leq p}$ 与 $\D^{\geq p^*}$，因为 $p^{**}=p$。特别地，它交换 $p$-\term{perverse-sheaf}{反常层}与 $p^*$-\term{perverse-sheaf}{反常层}，并交换 ${}^pH^i$ 与 ${}^{p^*}H^{-i}$。对由\term{stratum}{层片}组成的局部闭子集的嵌入 $j$，它交换 ${}^pj_!$ 与 ${}^{p^*}j_*$、${}^pj^!$ 与 ${}^{p^*}j^*$，并交换 ${}^pj_{!*}$ 与 ${}^{p^*}j_{!*}$。

在 $\D_c(X,R)$ 中，定义 $p$-\term{perverse-sheaf}{反常层}的条件可以重写为：对任意\term{stratum}{层片} $j\colon S\to X$，有
\[
 H^ij^*K=0\quad(i>p(S)),\qquad
 H^ij^*\mathbb DK=0\quad(i>p^*(S)).
\]
若所有\term{stratum}{层片}的维数均为偶数，则存在一个\term{self-dual}{自对偶}\term{perversity}{反常度}，由下式给出：
\[
 p_{1/2}(S)=-\tfrac12\dim S.
\]

\par\medskip\noindent\textbf{命题2.1.17.} 在2.1.16的假设下，若所有\term{stratum}{层片}均为偶数维，取\term{self-dual}{自对偶}\term{perversity}{反常度} $p_{1/2}$。设 $j\colon U\to X$ 是由\term{stratum}{层片}组成的开集，$A$ 是 $U$ 上的\term{self-dual}{自对偶}\term{perverse-sheaf}{反常层}。则 $j_{!*}A$ 是 $A$ 在 $\D_c(X,R)$ 中唯一满足下述条件的\term{self-dual}{自对偶}延拓 $P$：对每个\term{stratum}{层片} $S\subset X-U$，当 $i\geq-\frac12\dim S$ 时，$H^iP$ 在 $S$ 上为零。

$j_{!*}A$ 的\term{self-dual}{自对偶}性由 $A$ 的\term{self-dual}{自对偶}性和 $j_{!*}$ 的\term{self-dual}{自对偶}性得到。注意到这一点后，命题便是2.1.9与2.1.16的直接应用。

\SourcePage{71}
\par\medskip\noindent\textbf{注2.1.18.} 若 $U$ 可定向且\term{smooth}{光滑}，并\term{pure}{纯}为 $d$ 维，可以取 $A$ 为放在次数 $-d/2$ 的\term{constant-sheaf}{常值层} $R$。此时 $j_{!*}A$ 就是引言中记作 $\IC^\bullet$ 的“\term{intersection-complex}{交叉复形}”，而2.1.17就是Verdier对 $\IC^\bullet$ 的刻画。

\par\medskip\noindent\textbf{2.1.19.} 设 $(X,\mathcal S,\mathcal O)$ 如2.1.1，$\mathcal O^{\mathrm o}$ 为 $\mathcal O$ 的反环\term{sheaf}{层}，$p,q$ 为两个\term{perversity}{反常度}。下一个命题说明函子
\[
 \begin{aligned}
 \mathbin{\overset{\mathrm L}{\otimes}}\colon&\D^-(X,\mathcal O^{\mathrm o})\times\D^-(X,\mathcal O)\longrightarrow\D^-(X,\Z),\\
 \RHom\colon&\D^-(X,\mathcal O)\times\D^+(X,\mathcal O)\longrightarrow\D^+(X,\Z)
 \end{aligned}
\]
与\term{perversity}{反常度}的关系。若 $\mathcal O$ 交换，当然可以将 $\Z$ 换成 $\mathcal O$。

\par\medskip\noindent\textbf{命题2.1.20.} 函子 $\overset{\mathrm L}{\otimes}$ 将 $\D^{\leq p}\times\D^{\leq q}$ 映入 $\D^{\leq p+q}$；函子 $\RHom$ 将 $\D^{\leq p}\times\D^{\geq q}$ 映入 $\D^{\geq(q-p)}$。

第一个断言显然，因为 $\overset{\mathrm L}{\otimes}$ 可以逐点计算；写出它只是为了对称。设 $K\in\D^{\leq p}$、$L\in\D^{\geq q}$。分解归约可使我们假定 $K=j_!A[-n]$，其中 $A$ 是\term{stratum}{层片} $S$ 上的\term{sheaf}{层}，$j$ 是 $S$ 到 $X$ 的嵌入，且 $n\leq p(S)$。于是
\[
 \RHom(j_!A[-n],L)=j_*\RHom(A[-n],j^!L).
\]
按假设，$j^!L\in\D^{\geq q(S)}(S,\mathcal O)$，右端的 $\RHom$ 属于
\[
 \D^{\geq q(S)-n}(S,\Z)\subset\D^{\geq q(S)-p(S)}(S,\Z).
\]
而 $j_*$ 对任意\term{perversity}{反常度}都左 \term{t-exact}{$t$-正合}，故 $\RHom(j_!A[-n],L)\in\D^{\geq q-p}$。

令 $p=q$，得到：

\par\medskip\noindent\textbf{推论2.1.21.} 对 $K\in\D^{\leq p}(X,\mathcal O)$ 与 $L\in\D^{\geq p}(X,\mathcal O)$，有 $H^i\RHom(K,L)=0$（$i<0$）。

此推论也可由 $\Hom(K,L[i])=0$（$i<0$，见2.1.4）经\term{localization}{局部化}得到。

\par\medskip\noindent\textbf{推论2.1.22.} 对 $K\in\D^{\leq p}(X,\mathcal O)$ 与 $L\in\D^{\geq p}(X,\mathcal O)$，预\term{sheaf}{层}\[U\longmapsto\Hom_{\D(U,\mathcal O)}(K|_U,L|_U)\]是\term{sheaf}{层}。

事实上，由于 $H^i\RHom(K,L)$ 在 $i<0$ 时为零，局部到整体的\term{spectral-sequence}{谱序列}给出
\[
 \Hom_{\D(U,\mathcal O)}(K|_U,L|_U)=H^0(U,H^0\RHom(K,L)).
\]

\SourcePage{72}
在 $X$ 的每个开集 $U$ 上，赋予由 $\mathcal S$ 诱导的\term{stratification}{分层}和由 $p$ 诱导的\term{perversity}{反常度}。于是，$X$ 上的 $p$-\term{perverse-sheaf}{反常层}限制到 $U$ 上仍为 $p$-\term{perverse-sheaf}{反常层}。有：

\par\medskip\noindent\textbf{推论2.1.23.} $X$ 的各开集上的 $p$-\term{perverse-sheaf}{反常层}组成一个\term{stack}{叠}。

\emph{证明}（使用3.2；2.2.19中还给出了一个不依赖3.2的变式的证明）。由2.1.22，$p$-\term{perverse-sheaf}{反常层}组成一个预\term{stack}{叠}，即态射可以\term{recollement}{粘合}。对象的\term{recollement}{粘合}则应用3.2.4；由2.1.21，其假设得到满足。

\par\medskip\noindent\textbf{2.1.24. 变式。} 除了取 $X$ 为\term{stratification}{分层}拓扑空间，也可以取 $X$ 为\term{stratification}{分层}的\term{topos}{拓扑斯}。一个重要情形是：$X$ 为带\term{stratification}{分层}的\term{scheme}{概形} $S$ 的\term{etale}{平展}\term{topos}{拓扑斯}。

\SourcePage{73}
\section*{2.2. 概形}
\addcontentsline{toc}{section}{2.2. 概形}

\par\medskip\noindent\textbf{2.2.1.} 先以复代数簇为模型来讨论。将2.1应用于下述情形。

\emph{空间：} 复数域 $\C$ 上的分离\term{finite-type}{有限型}\term{scheme}{概形}的复点集 $X(\C)$，赋予通常拓扑。

\emph{\term{perversity}{反常度}：} 只考虑\term{stratum}{层片}的\term{perversity}{反常度} $p(S)$ 仅依赖其维数的情形。这里的维数总是偶数。因此，需要给定从非负偶数到 $\Z$ 的函数，仍记为 $p$。置 $p^*(n)=-n-p(n)$，这是\term{duality}{对偶}\term{perversity}{反常度}函数。假定 $p$ 与 $p^*$ 都递减，即
\[
 n\leq m\quad\Longrightarrow\quad0\leq p(n)-p(m)\leq m-n.
\]

\emph{\term{stratification}{分层}：} 回顾，每个 $X$ 的代数\term{stratification}{分层}都有一个Whitney细化（见 J.-L. Verdier，\emph{Stratifications de Whitney et théorème de Bertini-Sard}〔Whitney\term{stratification}{分层}与Bertini--Sard定理〕，Invent. Math. \textbf{36} (1976), 295--312）。所用\term{stratification}{分层}为 $X$ 的代数Whitney\term{stratification}{分层}，各\term{stratum}{层片}等维。它们满足2.1.13的(a)、(b)、(c)；我们只需要这些条件。下面在 $\C$ 上的讨论中，将它们简称为\term{stratification}{分层}。对每个\term{stratification}{分层} $\mathcal S$，其上的\term{perversity}{反常度} $p_{\mathcal S}$ 定义为
\[
 p_{\mathcal S}(S)=p(2\dim_{\mathrm{alg}}S)=p(\dim_{\mathrm{top}}S).
\]
以后以 $\dim S$ 简记代数维数 $\dim_{\mathrm{alg}}S$。

设 $R$ 为左\term{noetherian}{诺特}环。\emph{\term{constructible}{可构造}\term{derived-category}{导出范畴}} $\D_c(X(\C),R)$ 是 $\D(X(\C),R)$ 中由\term{cohomology}{上同调}\term{sheaf}{层} $H^iK$ 都\term{constructible}{可构造}的\term{complex}{复形} $K$ 组成的全子范畴。范畴
\[
 \D_c^b(X(\C),R)=\D^b(X(\C),R)\cap\D_c(X(\C),R)
\]
是 $\D(X(\C),R)$ 的各子范畴 $\D_{\mathcal S}^b(X(\C),R)$ 的\term{filtration}{滤过}并。在每个 $\D_{\mathcal S}(X(\C),R)$ 上，\term{perversity}{反常度} $p_{\mathcal S}$ 定义一个 \term{t-structure}{$t$-结构}。由2.1.14及对 $p$ 所作的假设，随着 $\mathcal S$ 逐渐变细，这些 \term{t-structure}{$t$-结构}相互诱导。取极限便得到 $\D_c^b(X(\C),R)$ 上\emph{\term{perversity}{反常度}为 $p$ 的 \term{t-structure}{$t$-结构}}
\[
 (\D_c^{b,\leq p}(X(\C),R),\D_c^{b,\geq p}(X(\C),R)).
\]
仿照2.1.4，可以由此得到 $\D_c(X(\C),R)$ 上的一个 \term{t-structure}{$t$-结构}：若对 $0\leq i\leq\dim X$，$a\leq p(2i)\leq b$，则定义 $\D_c^{\leq p}(X(\C),R)$（分别地，$\D_c^{\geq p}(X(\C),R)$）为
\SourcePage{74}
$\D_c(X(\C),R)$ 中满足以下条件的 $K$ 所成的子范畴：$i>b$（分别地，$i<a$）时 $H^iK=0$，且 $\tau_{[a,b]}K\in\D_c^{b,\leq p}(X(\C),R)$（分别地，$\D_c^{b,\geq p}(X(\C),R)$）。

对 $S\subset X$，置 $p(S)=p(2\dim S)$。

\par\medskip\noindent\textbf{命题2.2.2.} 设 $K\in\D_c(X(\C),R)$。下列条件\term{equivalence}{等价}。

(i) $K\in\D_c^{\leq p}(X(\C),R)$（分别地，$\D_c^{\geq p}(X(\C),R)$）。

(ii) $X$ 的每个不可约子簇 $S'$ 都有一个Zariski稠密开集 $S$，使得记 $S(\C)$ 到 $X(\C)$ 的嵌入为 $i$ 后，有 $H^ii^*K=0$（$i>p(S)$）（分别地，$H^ii^!K=0$，$i<p(S)$）。

若 $\mathcal T$ 是有限个Zariski局部闭、\term{smooth}{光滑}等维子集组成的族，其并为 $X$，且对 $S\in\mathcal T$，记其嵌入为 $i_S$ 后，$H^ii_S^*K$ 与 $H^ii_S^!K$ 均局部常值，则上述条件还\term{equivalence}{等价}于：

(iii) 对每个 $S\in\mathcal T$，$H^ii_S^*K=0$（$i>p(S)$）（分别地，$H^ii_S^!K=0$，$i<p(S)$）。

先设 $K\in\D_c^b$。对任意充分细的\term{stratification}{分层} $\mathcal S$，均有 $K\in\D_{\mathcal S}^b$。特别地，存在满足所述性质的 $\mathcal T$，因此只须证明 (i)$\Rightarrow$(ii)$\Rightarrow$(iii)$\Rightarrow$(i)。(ii)$\Rightarrow$(iii) 显然。验证 (i)$\Rightarrow$(ii) 时，只须取充分细的 $\mathcal S$，使 $K\in\D_{\mathcal S}^b$，从而 $K\in\D_{\mathcal S}^{b,\leq p}$（分别地，$\D_{\mathcal S}^{b,\geq p}$），并使 $S'$ 为某个\term{stratum}{层片}的闭包。再证 (iii)$\Rightarrow$(i)：取比 $\mathcal T$ 更细的 $\mathcal S$，使 $K\in\D_{\mathcal S}^b$。由对 $p$ 的假设以及2.1.4的证明，$K\in\D_{\mathcal S}^{b,\leq p}$（分别地，$\D_{\mathcal S}^{b,\geq p}$），于是得到(i)。

一般情形中，取 $a,b$，使对 $0\leq i\leq\dim X$ 都有 $a\leq p(2i)\leq b$。在(i)、(ii)或(iii)中把 $K$ 换成 $\tau_{\geq a}K$（分别地，$\tau_{\leq b}K$），所得条件与原条件\term{equivalence}{等价}。因此只须证明每个条件都蕴含 $H^iK=0$（$i>b$；分别地，$i<a$）。对(i)，这已包含在定义中。对(ii)和(iii)，重复2.1.2之后(2.1.2.1)的证明即可。对(iii)，先将 $\mathcal T$ 换成一个更细的\term{stratification}{分层}。对(ii)，对 $\dim S'$ 作下降归纳，证明：对每个不可约子簇 $S'$ 和每个 $i>b$（分别地，$i<a$），$S'$ 都有一个稠密开集 $S$，使 $H^iK$ 在 $S$ 上为零。

\par\medskip\noindent\textbf{2.2.3.} 若 $U$ 是 $X$ 的一个局部闭子集，则对每个充分细的\term{stratification}{分层}，$U$ 都是\term{stratum}{层片}的并。因此，可以
\SourcePage{75}
将2.1.6（$j_!,j^!,j^*,j_*$ 的正合性质）应用于 $U$ 到 $X$ 的嵌入 $j$。使用2.1.7的记号 $\tau_{\leq p},H^p,{}^pj_!,\ldots$；其中的结论仍然成立。对 $X$ 的Zariski开集 $U$ 及其补闭集 $F$，$\D_c(X(\C),R)$ 上\term{perversity}{反常度}为 $p$ 的 \term{t-structure}{$t$-结构}由 $\D_c(U(\C),R)$ 与 $\D_c(F(\C),R)$ 上的结构\term{recollement}{粘合}而来。因此，如2.1.8中一样，可以应用1.4的形式体系。

特别地，对 $U$ 上的 $p$-\term{perverse-sheaf}{反常层} $A$，由1.4.23，
\begin{equation*}
 j_{!*}A={}^p\tau^F_{\leq-1}j_*A.\tag{2.2.3.1}
\end{equation*}
假定 $\dim F\leq d$，置 $t=p(2d)$。由于 $\tau_{<t}i^*j_*A\in{}^p\D_c^{<0}$，且三角 $(\tau_{<t}i^*j_*A,\allowbreak A,\allowbreak\tau_{\geq t}i^*j_*A)$ 是\term{distinguished-triangle}{区别三角}\footnote{译注：本三角的中间对象照原文记为 $A$；原书勘误第27条将它更正为 $i^*j_*A$。}，若 $\tau_{\geq t}i^*j_*A\in{}^p\D_c^{\geq0}$，则有
\[
 {}^p\tau_{<0}i^*j_*A\xrightarrow{\sim}\tau_{<t}i^*j_*A,
 \qquad\text{从而}\qquad
 {}^p\tau^F_{<0}j_*A\xrightarrow{\sim}\tau^F_{<t}j_*A.
\]
特别地，将2.2.2(iii)应用于 $F$ 和 $\mathcal T=\{F\}$，得到：

\par\medskip\noindent\textbf{命题2.2.4.} 若 $F$ \term{smooth}{光滑}且维数为 $d$，并且 $H^ii^*j_*A$ 对 $i\geq t:=p(2d)$ 都在 $F$ 上局部常值，则
\begin{equation*}
 j_{!*}A=\tau^F_{\leq t-1}j_*A.\tag{2.2.4.1}
\end{equation*}
此式结合 $j_{!*}$ 的传递性，是2.1.11的对应结果。同样，在对 $F$ 的相同假设下，若 $H^ii^*j_*A$ 的局部常值性对 $i\geq t-1$（分别地，$i\geq t+1$）成立，则由1.4.23分别有
\begin{equation*}
 {}^pj_!A=\tau^F_{\leq t-2}j_*A,\tag{2.2.4.2}
\end{equation*}
\begin{equation*}
 {}^pj_*A=\tau^F_{\leq t}j_*A.\tag{2.2.4.3}
\end{equation*}

\par\medskip\noindent\textbf{命题2.2.5.} 若 $f\colon X\to Y$ 为拟有限态射，则函子 $f_!$ 与 $f^*$ 为右 \term{t-exact}{$t$-正合}，它们的右\term{adjunction}{伴随} $f^!$ 与 $f_*$ 为左 \term{t-exact}{$t$-正合}。

对 $K\in\D_c$，条件 $K\in{}^p\D_c^{\leq0}$ \term{equivalence}{等价}于对各 $H^iK$ 的\term{support}{支撑}维数的限制
\[
 p(2\dim\Supp H^iK)\geq i.
\]
由于函子 ${}^0f_!$ 正合，$H^if_!K={}^0f_!H^iK$，从而 $\Supp H^if_!K=f(\Supp H^iK)^-$，且 $\dim\Supp H^if_!K=\dim\Supp H^iK$。同样，$H^if^*K={}^0f^*H^iK$，$\Supp H^if^*K=f^{-1}\Supp H^iK$，且 $\dim\Supp H^if^*K\leq\dim\Supp H^iK$。这证明了 $f_!$ 与 $f^*$ 的右 \term{t-exact}{$t$-正合}性。$f^!$ 与 $f_*$ 的左 \term{t-exact}{$t$-正合}性可由\term{adjunction}{伴随}性推出（1.3.17(iii)）。

\SourcePage{76}
\emph{注。} 这里关键在于，\term{perversity}{反常度}条件是根据\term{stratum}{层片}的维数，而非余维数定义的（当 $X,Y$ 维数不同时尤其如此）。

\par\medskip\noindent\textbf{推论2.2.6.} (i) 若 $f$ 有限，则 $f_!=f_*$ 为 \term{t-exact}{$t$-正合}。

(ii) 若 $f$ \term{etale}{平展}，则 $f^!=f^*$ 为 \term{t-exact}{$t$-正合}。

\par\medskip\noindent\textbf{2.2.7.} 对拟有限的 $f$ 以及 $X$ 上的 $p$-\term{perverse-sheaf}{反常层} $F$，$f_!F\in{}^p\D^{\leq0}$、$f_*F\in{}^p\D^{\geq0}$。因此，$f_!F$ 到 $f_*F$ 的自然态射有分解
\begin{equation*}
 f_!F\longrightarrow{}^pf_!F\longrightarrow{}^pf_*F\longrightarrow f_*F.\tag{2.2.7.1}
\end{equation*}
置 $f_{!*}F=\im({}^pf_!F\to{}^pf_*F)$。

\par\medskip\noindent\textbf{2.2.8.} 如2.1.20，若\term{perversity}{反常度} $p,q$ 对 $n\leq m$ 满足
\[
 (p+q)(n)-(p+q)(m)\leq m-n,
\]
则 $\overset{\mathrm L}{\otimes}$ 将 $\D_c^{-,\leq p}\times\D_c^{-,\leq q}$ 映入 $\D_c^{\leq p+q}$；若 $q-p$ 递减，则 $\RHom$ 将 $\D_c^{-,\leq p}\times\D_c^{+,\geq q}$ 映入 $\D_c^{+,\geq(q-p)}$。特别地，对 $K\in\D_c^{\leq p}$、$L\in\D_c^{\geq p}$，$H^i\RHom(K,L)=0$（$i<0$），且 $U\mapsto\Hom_{\D(U)}(K|_U,L|_U)$ 是\term{sheaf}{层}。如2.1.23，$p$-\term{perverse-sheaf}{反常层}组成一个\term{stack}{叠}。

若 $R$ 为交换域，Verdier\term{duality}{对偶}交换 $\D^{>p}$ 与 $\D^{<p^*}$、$\D^{\geq p}$ 与 $\D^{\leq p^*}$，以及 $p$-\term{perverse-sheaf}{反常层}与 $p^*$-\term{perverse-sheaf}{反常层}（参见2.1.16），并交换 ${}^pH^j$ 与 ${}^{p^*}H^{-j}$。

\par\medskip\noindent\textbf{2.2.9.} 设 $X$ 是域 $k$ 上的\term{finite-type}{有限型}\term{scheme}{概形}，$\ell$ 为与 $k$ 的特征互素的素数。在 $X$ 上只考虑\term{etale}{平展}拓扑。我们的目标是定义 $X$ 上的 $\Q_\ell$-\term{perverse-sheaf}{反常层}。$\Z/\ell\Z$-模\term{sheaf}{层}的情形稍为简单，先从它开始。

\par\medskip\noindent\textbf{2.2.10.} 由于野分歧现象的存在，我们不再有能够起到Whitney\term{stratification}{分层}作用的\term{stratification}{分层}。需要考虑二元组 $(\mathcal S,\mathcal L)$，其中：

(a) $\mathcal S$ 把 $X$ 划分为有限个局部闭子集（称为\term{stratum}{层片}）。每个\term{stratum}{层片}的闭包是若干\term{stratum}{层片}的并；且在 $\bar k$ 上，每个\term{stratum}{层片}的约化\term{scheme}{概形}都\term{smooth}{光滑}、各处具有相同维数。

(b) $\mathcal L$ 对每个\term{stratum}{层片} $S$ 给定一个有限集合 $\mathcal L(S)$，其元素为 $S$ 上局部常值 $\Z/\ell\Z$-模\term{sheaf}{层}的同构类，
\SourcePage{77}
并且在所有局部常值 $\Z/\ell\Z$-模\term{sheaf}{层}组成的范畴中不可约。

若一个 $\Z/\ell\Z$-模\term{sheaf}{层}在每个\term{stratum}{层片} $S$ 上的限制局部常值，且为同构类属于 $\mathcal L(S)$ 的\term{sheaf}{层}的有限次迭代\term{extension}{扩张}，则称它为\emph{$(\mathcal S,\mathcal L)$-\term{constructible}{可构造}\term{sheaf}{层}}。用 $\D_{\mathcal S,\mathcal L}^b(X,\Z/\ell\Z)$ 表示 $\D^b(X,\Z/\ell\Z)$ 中由 $(\mathcal S,\mathcal L)$-\term{constructible}{可构造}\term{complex}{复形}组成的全子范畴，即要求各 $H^iK$ 都是 $(\mathcal S,\mathcal L)$-\term{constructible}{可构造}的。若 $\mathcal S$ 的每个\term{stratum}{层片} $S$ 都是 $\mathcal S'$ 的\term{stratum}{层片}的并，且 $\mathcal L(S)$ 中的各 $F$ 都是 $(\mathcal S',\mathcal L')$-\term{constructible}{可构造}的，即为 $(\mathcal S'|_S,\mathcal L'|_{(\mathcal S'|_S)})$-\term{constructible}{可构造}的，则称 $(\mathcal S',\mathcal L')$ \emph{细化} $(\mathcal S,\mathcal L)$。

最后再要求：

(c) 对 $S\in\mathcal S$、$F\in\mathcal L(S)$，若 $j$ 是 $S$ 到 $X$ 的嵌入，则各 $R^nj_*F$ 都是 $(\mathcal S,\mathcal L)$-\term{constructible}{可构造}的。

和2.1.13中一样，这个条件保证：对于局部闭且由\term{stratum}{层片}组成的 $U\xhookrightarrow{j}V\subset X$，函子 $j_*,j_!,j^*,j^!$ 将 $\D_{\mathcal S,\mathcal L}^b(U,\Z/\ell\Z)$ 与 $\D_{\mathcal S,\mathcal L}^b(V,\Z/\ell\Z)$ 相互映入。

由 $Rj_*$ 的\term{constructible}{可构造}性定理（[3] 1.1），对任意满足(a)、(b)的系统 $(\mathcal S',\mathcal L')$，存在一个满足(a)、(b)、(c)的细化 $(\mathcal S,\mathcal L)$。事实上，若将 $X$ 换成所有维数 $\geq n$ 的\term{stratum}{层片}之并时，(a)、(b)、(c)已成立，则\term{constructible}{可构造}性定理使我们能够将 $(\mathcal S,\mathcal L)$ 细化为 $(\mathcal S',\mathcal L')$，不改变维数 $\geq n$ 的\term{stratum}{层片}，并使(a)、(b)、(c)在将 $X$ 换成维数 $\geq n-1$ 的\term{stratum}{层片}之并时成立。迭代这一构造即可。

\par\medskip\noindent\textbf{2.2.11.} 设 $p$ 如2.2.1，给出\term{perversity}{反常度}函数
\[
 p_{\mathcal S}\colon\mathcal S\to\Z,\qquad S\mapsto p(2\dim S).
\]
按前述方法，由 $p$ 得到 $\D_{\mathcal S,\mathcal L}^b(X,\Z/\ell\Z)$ 上的一个 \term{t-structure}{$t$-结构}。随着 $(\mathcal S,\mathcal L)$ 变细，这些 \term{t-structure}{$t$-结构}相互诱导：SGA 4 XVI 3.7的\term{pure}{纯}性定理给出2.1.14的对应结果。取极限，在各 $\D_{\mathcal S,\mathcal L}^b(X,\Z/\ell\Z)$ 的\term{filtration}{滤过}并 $\D_c^b(X,\Z/\ell\Z)$ 上得到一个 \term{t-structure}{$t$-结构}
\[
 (\D_c^{b,\leq p}(X,\Z/\ell\Z),\D_c^{b,\geq p}(X,\Z/\ell\Z)).
\]
如2.2.1，可由此得到 $\D_c(X,\Z/\ell\Z)$ 上的一个 \term{t-structure}{$t$-结构}。2.2.2至2.2.8仍成立，证明基本相同。只需将Verdier\term{duality}{对偶}的形式体系换成\term{etale}{平展}\term{cohomology}{上同调}的\term{duality}{对偶}形式体系（SGA 4 XVIII）。

\par\medskip\noindent\textbf{2.2.12.} 对 $X$ 的一点 $x$（不要求闭），将 $x$ 到 $X$ 的嵌入 $i_x$ 分解为 $x\xrightarrow{j}\overline{\{x\}}\xrightarrow{i}X$，定义 $i_x^!:=j^*i^!$。用
\SourcePage{78}
这个记号，2.2.2(ii)对 $\D_c^{\leq p}$（分别地，$\D_c^{\geq p}$）的刻画可以更简洁地表述为：

(ii)$^*$ 对 $X$ 的每个点 $x$（不要求闭），记 $\dim x=\dim\overline{\{x\}}$，则 $H^ii_x^*K=0$（$i<p(2\dim x)$）（分别地，$H^ii_x^!K=0$，$i>p(2\dim x)$）。\footnote{译注：此处两个不等号照原文。原书勘误第31、32条分别将它们更正为 $i>p(2\dim x)$ 与 $i<p(2\dim x)$。}

\par\medskip\noindent\textbf{2.2.13.} 对 $\Z/\ell^n\Z$-模\term{sheaf}{层}（更一般地，对 $R$-模\term{sheaf}{层}，其中 $R$ 为有限的 $\Z/\ell^n\Z$-代数），采用同样的方法。方便的做法是仍令 $\mathcal L$ 对每个\term{stratum}{层片}给定有限个不可约局部常值 $\Z/\ell\Z$-模\term{sheaf}{层}；若 $K\in\D_c^b(X,R)$ 的各 $\ell^iH^*K/\ell^{i+1}H^*K$ 作为 $\Z/\ell\Z$-模\term{sheaf}{层}都是 $(\mathcal S,\mathcal L)$-\term{constructible}{可构造}的，就称 $K$ 为 $(\mathcal S,\mathcal L)$-\term{constructible}{可构造}的。

\par\medskip\noindent\textbf{2.2.14.} 对 $\Z_\ell$-\term{sheaf}{层}（分别地，$\Q_\ell$-\term{sheaf}{层}等），也采用同样的方法；前提是先定义\term{triangulated-category}{三角范畴} $\D_c^b(X,\Z_\ell)$（分别地，$\D_c^b(X,\Q_\ell)$ 等），使它满足通常的变差形式体系，并定义其“自然” \term{t-structure}{$t$-结构}，要求其\term{heart}{心}为 $X$ 上\term{constructible}{可构造}的 $\Z_\ell$-\term{sheaf}{层}（分别地，$\Q_\ell$-\term{sheaf}{层}等）组成的\term{abelian-category}{阿贝尔范畴}（SGA 5 VII 1.1）。基域为有限域或代数闭域时，或更一般地，当 $k$ 的每个有限\term{extension}{扩张} $k'$ 的群 $H^i(\Gal(\bar k/k'),\Z/\ell\Z)$ 都有限时，[1] 1.1.2提出了一个办法：置
\begin{equation*}
 \D_c^b(X,\Z_\ell):=2\!\varprojlim_n\D_{\mathrm{ctf}}^b(X,\Z/\ell^n\Z).\tag{2.2.14.1}
\end{equation*}
其中ctf表示“\term{constructible}{可构造}且Tor维数有限”。过渡函子取为标量\term{extension}{扩张}函子
\[
 -\mathbin{\overset{\mathrm L}{\otimes}}_{\Z/\ell^n\Z}\Z/\ell^m\Z.
\]
为使它们将 $\D^b$ 映入 $\D^b$，必须限制于Tor维数有限的对象。

下面的命题结合[3]中的有限性定理，保证得到的确为一个\term{triangulated-category}{三角范畴}。

\par\medskip\noindent\textbf{命题2.2.15.} 设 $(\mathcal D_n)_{n\in\mathbb N}$ 是\term{triangulated-category}{三角范畴}的逆系统，过渡函子 $T_{m,n}\colon\mathcal D_n\to\mathcal D_m$ 为\term{exact-functor}{正合函子}。若对每个 $n\in\mathbb N$ 及 $K,L\in\mathcal D_n$，$\Hom(K,L)$ 都有限，则各 $\mathcal D_n$ 的2-\term{projective-limit}{逆极限} $\mathcal D$ 配上如下三角族后为\term{triangulated-category}{三角范畴}：三角 $T$ 在每个 $\mathcal D_n$ 中的像均为\term{distinguished-triangle}{区别三角}。

设 $u\colon X\to Y$ 是 $\mathcal D$ 中的态射，其在 $\mathcal D_n$ 中的像为 $u_n\colon X_n\to Y_n$。为
\SourcePage{79}
验证 $u$ 为某个\term{distinguished-triangle}{区别三角} $T$ 的底边（TR1），对每个 $n$ 取一个底边为 $u_n$ 的\term{distinguished-triangle}{区别三角} $T_n$。在各 $\mathcal D_n$ 中应用TR3，构造底边为 $u_n$ 的恒等自同构的态射 $T_{n,n+1}(T_{n+1})\to T_n$。它们都是同构。因此可以把所求的 $T$ 定义为各 $T_n$ 的\term{projective-limit}{逆极限}。

再设 $T,T'$ 是 $\mathcal D$ 中的两个\term{distinguished-triangle}{区别三角}，$f\colon\operatorname{base}(T)\to\operatorname{base}(T')$：
\[
 \begin{tikzcd}[ampersand replacement=\&,column sep=large]
 X\arrow[r,"{u}"]\arrow[d,"{f_X}"']\&Y\arrow[r]\arrow[d,"{f_Y}"]\&Z\arrow[r,"{[1]}"]\&{}\\
 X'\arrow[r]\&Y'\arrow[r]\&Z'\arrow[r,"{[1]}"]\&{}
 \end{tikzcd}
\]
对每个 $n$，设 $T_n,T'_n,f_n$ 是 $T,T',f$ 在 $\mathcal D_n$ 中的像，$H_n$ 是诱导 $f_n$ 的三角态射 $T_n\to T'_n$ 的集合。函子 $T_{n,n+1}$ 把 $H_{n+1}$ 映入 $H_n$，各 $H_n$ 的\term{projective-limit}{逆极限} $H$ 就是诱导 $f$ 的 $k\colon T\to T'$ 所成的集合。按假设，每个 $H_n$ 有限；由 $\mathcal D_n$ 中的TR3，它又非空。因此 $H\ne\varnothing$，这就验证了 $\mathcal D$ 中的TR3。$\mathcal D$ 中TR4的验证相同，TR2的验证则显然。

\par\medskip\noindent\textbf{2.2.16.} 定义(2.2.14.1)的优点在于，$\D_c^b(X,\Z_\ell)$ 的变差形式体系只须取极限便可得到。自然 \term{t-structure}{$t$-结构}的定义则较微妙，因为 $\D_c^b(X,\Z/\ell^n\Z)$ 中的 $\tau_{\leq i}$ 不保持Tor维数有限的条件，且不与标量\term{extension}{扩张}交换。定义从 $\D_c^b(X,\Z_\ell)$ 到\term{constructible}{可构造} $\Z_\ell$-\term{sheaf}{层}的\term{abelian-category}{阿贝尔范畴}的函子 $H^i$ 如下：对由系统 $K_n\in\D_c^b(X,\Z/\ell^n\Z)$ 给出的 $K$，取各 $H^i(K_n)$ 的\term{projective-limit}{逆极限}，作为对应的 $\Z_\ell$-\term{sheaf}{层}。将 $K_n$ 记为 $K\otimes_{\Z_\ell}\Z/\ell^n\Z$。有通常的\term{exact-sequence}{正合列}
\begin{equation*}
\begin{aligned}
 0\longrightarrow H^i(K)\otimes_{\Z_\ell}\Z/\ell^n\Z
 &\longrightarrow H^i(K\otimes_{\Z_\ell}\Z/\ell^n\Z)\\
 &\longrightarrow\operatorname{Tor}_1(H^{i+1}(K),\Z/\ell^n\Z)\longrightarrow0.
\end{aligned}\tag{2.2.16.1}
\end{equation*}
随着 $n$ 变化，它们构成一个逆系统；各 $\operatorname{Tor}_1$ 组成的逆系统本质为零。以 $H^iK=0$（$i>0$；分别地，$i<0$）定义 $\D_c^b(X,\Z_\ell)^{\leq0}$（分别地，$\D_c^b(X,\Z_\ell)^{\geq0}$）。要保证确实得到 \term{t-structure}{$t$-结构}，关键是构造算子 $\tau_{\leq i}$；[1] 1.1.2中完成了这一构造。函子 $H^0$ 诱导范畴\term{equivalence}{等价}
\[
 \bigl(\D_c^b(X,\Z_\ell)^{\leq0}\cap\D_c^b(X,\Z_\ell)^{\geq0}\bigr)
 \longrightarrow(\text{\term{constructible}{可构造}的 $\Z_\ell$-\term{sheaf}{层}}).
\]
\term{exact-sequence}{正合列}(2.2.16.1)表明：$K\in\D_c^b(X,\Z_\ell)^{\leq0}$ 当且仅当对某个（\term{equivalence}{等价}地，每个）$n$，$K\otimes_{\Z_\ell}\Z/\ell^n\Z$
\SourcePage{80}
属于 $\D_c^b(X,\Z/\ell^n\Z)^{\leq0}$。

O. Gabber 提出了一个良好且无条件的 $\D_c^b(X,\Z_\ell)$ 定义。将它给出的\term{triangulated-category}{三角范畴}记为 $\widehat\D_c(X,\Z_\ell)$\footnote{译注：此处记号照原文；原书勘误第33条指出，应补上上标 $b$，即 $\widehat\D_c^b(X,\Z_\ell)$。}。有一个\term{conservative}{保守函子}
\[
 \widehat\D_c^b(X,\Z_\ell)\longrightarrow\D_c^b(X,\Z_\ell),
\]
在同构类集合上给出双射。对 $\widehat\D_c^b(X,\Z_\ell)$ 中的 $K,L$，设模 $\ell^n$ 约化分别为 $K_n,L_n$，则有短\term{exact-sequence}{正合列}
\[
 0\longrightarrow\varprojlim{}^1\Hom(K_n,L_n)
 \longrightarrow\Hom(K,L)\longrightarrow\varprojlim\Hom(K_n,L_n)
 \longrightarrow0.
\]
\noindent\footnote{译注：本短\term{exact-sequence}{正合列}的首项照原文；原书勘误第34条指出，首项应为 $\varprojlim{}^1\Hom^{-1}(K_n,L_n)$，原式漏印了 $-1$。}
当 $X$ 为域 $k$ 的谱、$\bar k$ 为其可分闭包、$V$ 是带有连续 $\Gal(\bar k/k)$ 作用的有限生成 $\Z_\ell$-模，且 $K$（分别地，$L$）为集中在次数 $0$ 的 $\Z_\ell$（分别地，$V$）时，在 $\widehat\D_c^b$ 中计算的 $\Hom^i(\Z_\ell,V)$，就是由连续上链\term{complex}{复形}计算的 $H^i(\Gal(\bar k/k),V)$。由于他尚未写出这个定义，我们只得限制在2.2.15适用的这个较为狭窄的框架内。

\par\medskip\noindent\textbf{2.2.17.} 设 $\D_{\mathcal S,\mathcal L}^b(X,\Z_\ell)$ 是 $\D_c^b(X,\Z_\ell)$ 中由满足下列条件的 $K$ 组成的全子范畴；由(2.2.16.1)，这些条件\term{equivalence}{等价}：(a)对每个（分别地，某个）$n$，$K\otimes_{\Z_\ell}\Z/\ell^n\Z\in\D_{\mathcal S,\mathcal L}^b(X,\Z/\ell^n\Z)$；(b)各 $H^iK\otimes_{\Z_\ell}\Z/\ell\Z$ 都是 $(\mathcal S,\mathcal L)$-\term{constructible}{可构造}的。当 $(\mathcal S,\mathcal L)$ 变化时，$\D_c^b$ 是各 $\D_{\mathcal S,\mathcal L}^b$ 的\term{filtration}{滤过}并。如2.2.10，在各 $\D_{\mathcal S,\mathcal L}^b$ 上\term{recollement}{粘合}得到 \term{t-structure}{$t$-结构}，再对 $(\mathcal S,\mathcal L)$ 取极限，即定义 $\D_c^b$ 上\term{perversity}{反常度}为 $p$ 的 \term{t-structure}{$t$-结构}。

如2.2.16，$K\in{}^p\D_c^{\leq0}(X,\Z_\ell)$ 当且仅当其模 $\ell$ 约化属于 ${}^p\D_c^{\leq0}(X,\Z/\ell\Z)$。

\par\medskip\noindent\textbf{2.2.18.} 回顾，按定义，
\[
 \D_c^b(X,\Q_\ell):=\D_c^b(X,\Z_\ell)\otimes_{\Z_\ell}\Q_\ell,
\]
而\term{constructible}{可构造} $\Q_\ell$-\term{sheaf}{层}的\term{abelian-category}{阿贝尔范畴}，是\term{constructible}{可构造} $\Z_\ell$-\term{sheaf}{层}范畴的 $\Q_\ell\otimes_{\Z_\ell}-$。全子范畴 $\D_{\mathcal S,\mathcal L}^b(X,\Q_\ell)$ 定义为 $\D_{\mathcal S,\mathcal L}^b(X,\Z_\ell)\otimes\Q_\ell$ 的像；它由各 $H^iK$ 都是某个 $(\mathcal S,\mathcal L)$-\term{constructible}{可构造} $\Z_\ell$-\term{sheaf}{层}的 $\Q_\ell\otimes-$ 的 $K$ 组成，这里的 $(\mathcal S,\mathcal L)$-\term{constructible}{可构造}指其模 $\ell$ 约化 $(\mathcal S,\mathcal L)$-\term{constructible}{可构造}。如前，通过\term{recollement}{粘合}以及对 $(\mathcal S,\mathcal L)$ 取极限，定义 $\D_c^b(X,\Q_\ell)$ 上\term{perversity}{反常度}为 $p$ 的 \term{t-structure}{$t$-结构}。对任意区间 $[a,b]$，自然函子
\[
 {}^p\D_c^{[a,b]}(X,\Z_\ell)\otimes\Q_\ell
 \longrightarrow{}^p\D_c^{[a,b]}(X,\Q_\ell)
\]
为\term{equivalence}{等价}。特别地（$a=b=0$），
\SourcePage{81}
$\Q_\ell$-\term{perverse-sheaf}{反常层}的\term{abelian-category}{阿贝尔范畴}由 $\Z_\ell$-\term{perverse-sheaf}{反常层}范畴经 $-\otimes_{\Z_\ell}\Q_\ell$ 得到。

设 $E_\lambda$ 是 $\Q_\ell$ 的一个有限\term{extension}{扩张}，$\mathcal O_\lambda$ 是其赋值环。可像处理 $\Z_\ell$-\term{sheaf}{层}一样处理 $\mathcal O_\lambda$-\term{sheaf}{层}，再经 $-\otimes_{\mathcal O_\lambda}E_\lambda$ 得到 $E_\lambda$-\term{sheaf}{层}。也可以由 $\Q_\ell$-\term{sheaf}{层}过渡到 $E_\lambda$-\term{sheaf}{层}：注意遗忘函子 $\omega$ 诱导一个\term{equivalence}{等价}
\[
 \D_c^b(X,E_\lambda)\longrightarrow
 \left\{\begin{array}{l}
 K\in\D_c^b(X,\Q_\ell),\text{ 配有一个}\\
 \Q_\ell\text{-代数态射 }E_\lambda\to\operatorname{End}(K)
 \end{array}\right\}.
\]
其\term{fully-faithful}{全忠实}性由\term{spectral-sequence}{谱序列}
\begin{equation*}
 E_2^{pq}=\Ext^p_{E_\lambda\otimes E_\lambda}
 \bigl(E_\lambda,\Hom^q(\omega K,\omega L)\bigr)
 \Longrightarrow\Hom(K,L)\tag{2.2.18.1}
\end{equation*}
得到，其中当 $p\ne0$ 时 $E_2^{pq}=0$（张量积 $E_\lambda\otimes E_\lambda$ 在 $\Q_\ell$ 上取，$\Hom^q$ 则在 $\D_c^b(X,\Q_\ell)$ 中计算）。本质满性来自：带有 $E_\lambda$ 作用的 $K$ 是 $K\otimes_{\Q_\ell}E_\lambda$ 的直和项。

\term{spectral-sequence}{谱序列}(2.2.18.1)由 $\D_c^b(X,\mathcal O_\lambda)$ 的一个类似\term{spectral-sequence}{谱序列}得到，后者通过取极限构造。若 $\mathcal O_\lambda$ 在 $\Z_\ell$ 上\term{etale}{平展}，则仍有 $E_2^{pq}=0$（$p\ne0$），一般情形则不成立。不过，这一\term{spectral-sequence}{谱序列}仍表明遗忘函子诱导一个\term{fully-faithful}{全忠实}函子
\[
 (\mathcal O_\lambda\text{-\term{perverse-sheaf}{反常层}})\longrightarrow
 (\text{带有 $\mathcal O_\lambda$ 作用的 $\Z_\ell$-\term{perverse-sheaf}{反常层}}).
\]
我们尚未验证它如预期那样是一个\term{equivalence}{等价}。

最后，对 $E_\lambda\subset\overline\Q_\ell$ 取归纳极限，得到 $\D_c^b(X,\overline\Q_\ell)$。

\par\medskip\noindent\textbf{2.2.19.} 在所有这些框架中，2.2.2至2.2.8和2.2.12仍然适用，证明基本相同。

不过，2.1.23中关于\term{perverse-sheaf}{反常层}组成\term{stack}{叠}的证明不能原封不动地使用，因为它所依赖的3.2是在某个\term{sheaf}{层}范畴的\term{derived-category}{导出范畴}内写成的。下面给出另一种证明。如2.1.23，我们已经知道态射可以\term{recollement}{粘合}。需要证明的是：若 $(U_i)_{i\in I}$ 是 $X$ 的一个有限\term{etale}{平展}覆盖，则 $U_i$ 上任意带\term{recollement}{粘合}数据的\term{perverse-sheaf}{反常层}族 $A_i$
\SourcePage{82}
都来自 $X$ 上的某个\term{perverse-sheaf}{反常层} $A$。\term{recollement}{粘合}数据使我们能够对每个经过某个 $U_i$ 分解的\term{etale}{平展}态射 $j\colon V\to X$，定义 $V$ 上的\term{perverse-sheaf}{反常层} $A_V$；而且这一定义与 $V'\to V$ 的\term{inverse-image}{逆像}相容。

对任意态射 $f\colon Y\to X$ 和 $X$ 的任意不可约局部闭子簇 $i\colon S'\hookrightarrow X$，作笛卡尔图
\[
 \begin{tikzcd}[ampersand replacement=\&,column sep=large,row sep=large]
 T'\arrow[r,hook,"{i'}"]\arrow[d,"{f}"']\&Y\arrow[d,"{f}"]\\
 S'\arrow[r,hook,"{i}"']\&X
 \end{tikzcd}
\]
则在\term{derived-category}{导出范畴}中有 $i^!f_*=f_*i'^!$。若 $f$ 拟有限，则 $S'$ 有一个稠密开集 $S$，使 $T=f^{-1}S$ 在 $S$ 上有限，并且或为空，或\term{pure}{纯}为 $\dim S$ 维。因此，$f_*$ 将 $\D_c^{>p}(Y,\Q_\ell)$ 映入 $\D_c^{>p}(X,\Q_\ell)$，从而 ${}^pf_*:=H^0f_*$ 是从 $Y$ 上\term{perverse-sheaf}{反常层}到 $X$ 上\term{perverse-sheaf}{反常层}的一个左\term{exact-functor}{正合函子}，其左\term{adjunction}{伴随}为 ${}^pf^*:=H^0f^*$（1.3.17(iii)）。

设 $j_i$（分别地，$j_{ij}$）为 $U_i$（分别地，$U_{ij}:=U_i\times_XU_j$）到 $X$ 的投影。置
\[
 A:=\Ker\left(\prod_{i\in I}{}^pj_{i*}A_{U_i}
 \rightrightarrows\prod_{(i,j)\in I\times I}{}^pj_{ij*}A_{U_{ij}}\right).
\]
只需证明，由 $A\to j_{i*}A_{U_i}$ 经\term{adjunction}{伴随}得到的态射 ${}^pj_i^*A\to A_i=A_{U_i}$ 是同构。又只需在每次作\term{base-change}{基变换} $j_i$ 后证明。所考虑的\term{direct-image}{正像}和\term{inverse-image}{逆像}都与这种\term{localization}{局部化}交换，因此归结到某个 $U_j/X$ 有截面的情形。这时，各 $A_i$ 都是某个 $A_X$ 的\term{inverse-image}{逆像}，而 $A=A_X$ 可由通常的\term{homotopy}{同伦}论证得到：这一论证证明，两个互相细化的开覆盖给出相同的Čech\term{cohomology}{上同调}。
